% =====================================================================
\chapter{Existencia, unicidad, contraejemplos y redacción}
\label{tema:6}
% =====================================================================

Los temas anteriores se han centrado en enunciados universales. Este tema trata los enunciados existenciales y de unicidad, explica cómo se refuta un enunciado falso y termina con unas pautas para redactar demostraciones claras, que resumen todo el módulo. Se siguen los capítulos~7 y~9 y la §5.3 de~\cite{hammack}, la §3.6 de~\cite{velleman}, la §1.9 de~\cite{mcs}, la lista de comprobación de CS103~\cite{cs103} y el método de Polya~\cite{polya}.

\section{Demostraciones de existencia}

Para demostrar $\exists x\,P(x)$ basta encontrar \emph{un} elemento $a$, llamado \emph{testigo}, y comprobar $P(a)$. Una demostración de este tipo es \emph{constructiva}. A veces se prueba la existencia sin exhibir el testigo, por ejemplo por reducción al absurdo o examinando casos sin saber cuál se da; la demostración es entonces \emph{no constructiva}~\cite[§7.3–7.4]{hammack}.

\addtocounter{definicion}{1}%

\begin{proposicion}[Existencia no constructiva]\label{prop:irracional-potencia}
Existen números irracionales $a$ y $b$ tales que $a^{b}$ es racional.
\end{proposicion}

La prueba no dice cuál de los dos pares sirve. (Se sabe, por un teorema profundo de Gelfond y Schneider, que $c$ es irracional, pero la demostración anterior no lo necesita.)

\section{Demostraciones de unicidad}

Para demostrar $\exists!\,x\,P(x)$ se prueban dos cosas (\cref{def:unicidad}): la \emph{existencia}, $\exists x\,P(x)$, y la \emph{unicidad}, que se redacta así: «sean $x$ e $y$ tales que $P(x)$ y $P(y)$; entonces $x=y$». Una alternativa equivalente es suponer $x\neq y$ y llegar a una contradicción~\cite[§3.6]{velleman}.

\begin{proposicion}[Unicidad en la división euclídea]\label{prop:division-unicidad}
Sean $a\in\Z$ y $b\in\N$. Si $a=qb+r=q'b+r'$ con $q,q',r,r'\in\Z$ y $0\le r,r'<b$, entonces $q=q'$ y $r=r'$.
\end{proposicion}

Junto con la \cref{prop:division-existencia}, esto completa la demostración del \cref{teo:division}.

\addtocounter{definicion}{1}%

\section{Contraejemplos y refutaciones}

Para \emph{refutar} un enunciado se demuestra su negación~\cite[cap.~9]{hammack}. Según la forma del enunciado:
\begin{itemize}[leftmargin=1.4em]
  \item $\forall x\,P(x)$ se refuta con un \emph{contraejemplo}: un $a$ concreto con $\neg P(a)$.
  \item $\forall x\,(P(x)\Rightarrow Q(x))$ se refuta con un $a$ que cumple $P(a)$ y no cumple $Q(a)$.
  \item $\exists x\,P(x)$ se refuta demostrando $\forall x\,\neg P(x)$, lo que exige una demostración general, no un ejemplo.
\end{itemize}

\addtocounter{definicion}{1}%

\addtocounter{definicion}{1}%

\section{Cómo redactar una demostración}

Una demostración es un texto dirigido a un lector, no una lista de cálculos. Las pautas siguientes resumen las de la §5.3 de~\cite{hammack}, la §1.9 de~\cite{mcs} y la lista de comprobación de CS103~\cite{cs103}.

\begin{enumerate}[label=\arabic*.]
  \item \emph{Anunciar el plan.} Indicar qué se supone, qué se quiere demostrar y con qué técnica («Se procede por contrarrecíproco», «Por inducción sobre $n$»).
  \item \emph{Introducir cada variable.} Toda letra debe presentarse antes de usarse, con su universo: «Sea $n$ un entero impar»; «existe $k\in\Z$ con\dots». No reutilizar una letra para objetos distintos.
  \item \emph{Usar las definiciones, no repetirlas.} Desplegar la definición en el punto donde se necesita, aplicada a los objetos concretos.
  \item \emph{Escribir frases completas.} Las fórmulas forman parte de las frases. Los símbolos $\forall$, $\exists$, $\Rightarrow$ son para enunciados formales, no abreviaturas en la prosa.
  \item \emph{Que cada frase aporte.} Eliminar lo que no se usa después; justificar cada paso no evidente citando el resultado empleado.
  \item \emph{Terminar con la conclusión.} Decir explícitamente que se ha demostrado lo que se buscaba y cerrar con $\blacksquare$.
\end{enumerate}

Para \emph{encontrar} la demostración, antes de redactarla, sirve el método de cuatro fases de Polya~\cite{polya}: comprender el problema (datos, incógnita, condiciones; escribir el enunciado con cuantificadores), concebir un plan (¿se conoce un problema relacionado?, ¿sirve un caso particular o el contrarrecíproco?), ejecutarlo comprobando cada paso, y examinar la solución obtenida (¿se usaron todas las hipótesis?, ¿hay un argumento más directo?).

\addtocounter{definicion}{1}%

\addtocounter{definicion}{1}%

\addtocounter{definicion}{1}%

